\documentclass[11pt,a4paper]{article}
\usepackage[utf8]{inputenc}
\usepackage[T1]{fontenc}
\usepackage[ngerman]{babel}
\usepackage{amsmath,amssymb}
\usepackage{tikz}
\usetikzlibrary{trees,arrows.meta,positioning}
\usepackage[margin=2.5cm]{geometry}
\usepackage{enumitem}
\usepackage{booktabs}
\usepackage{algorithm}
\usepackage{algpseudocode}

\title{Greedy Set Cover\\[0.3em]\large Zweiphasiger Algorithmus zur precision-optimalen Regelkomposition}
\author{}
\date{}

\begin{document}
\maketitle

%% ---------------------------------------------------------------
\section{Problemstellung}

Gegeben sei ein Datensatz mit $n$ Zeilen, $m$ bin\"aren Pr\"adikaten $p_1, \dots, p_m$ (aufgeteilt in Feature-Gruppen) sowie einem bin\"aren Label~$\mathbf{y} \in \{0,1\}^n$ mit $\text{total\_pos} = \mathbf{1}^\top \mathbf{y}$ Positiven.

\medskip
\textbf{Ziel:} Finde eine \emph{Composition}
\[
  C = R_1 \;\lor\; R_2 \;\lor\; \cdots \;\lor\; R_r
\]
aus $r$ AND-Regeln, sodass die \textbf{Precision} maximiert wird, w\"ahrend ein gew\"unschter \textbf{Recall} erreicht wird. Im Unterschied zu den B\&B-Ans\"atzen wird kein fester min\_recall-Constraint gesetzt. Stattdessen liefert der Algorithmus \emph{alle Zwischenschritte} --- eine monoton steigende Recall-Folge --- aus der f\"ur jedes gew\"unschte Recall-Ziel die beste Composition abgelesen werden kann.

\medskip
\textbf{Parameter:}
\begin{itemize}[nosep]
  \item \texttt{max\_k}: Maximale Anzahl Pr\"adikate pro AND-Regel
  \item \texttt{min\_tp}: Minimale True Positives f\"ur Kandidatenregeln (Phase~1)
  \item \texttt{min\_new\_tp}: Minimale \emph{neue} True Positives, damit eine Regel in einem Schritt ausgew\"ahlt wird (Phase~2)
  \item \texttt{max\_candidates}: Obergrenze f\"ur die Kandidatenanzahl (Diversity-Filter)
\end{itemize}

%% ---------------------------------------------------------------
\section{Phase 1: Kandidatengenerierung (Apriori)}

Die Kandidatengenerierung nutzt denselben Apriori-Algorithmus wie der Apriori+B\&B-Ansatz. Die Anti-Monotonie der AND-Verkn\"upfung dient als Pruning-Kriterium:
\[
  \mathit{TP}(R \land p) \;\leq\; \mathit{TP}(R)
\]

\subsection{Level-weiser Aufbau}

\textbf{Level 1:} Erzeuge alle Ein-Pr\"adikat-Regeln $\{p_i\}$ mit $\mathit{TP}(p_i) \geq \text{min\_tp}$.

\medskip
\textbf{Level $\ell \to \ell+1$} (bis $\ell = \texttt{max\_k}$): F\"ur jede Regel $R$ aus Level~$\ell$ und jedes Pr\"adikat $p_j$ mit Gruppenindex $>$ h\"ochster Gruppenindex in $R$:
\begin{itemize}[nosep]
  \item Bilde $R' = R \land p_j$
  \item Falls $\mathit{TP}(R') \geq \text{min\_tp}$: behalte $R'$
  \item Sonst: verwerfe (Pruning)
\end{itemize}

\noindent Die Constraint \emph{maximal ein Pr\"adikat pro Feature-Gruppe} wird durch den aufsteigenden Gruppenindex sichergestellt.

\subsection{Scoring und Diversity-Filter}

Jede Kandidatenregel $R_i$ wird durch ihren Maskenvektor $\mathbf{m}_i \in \{0,1\}^n$ repr\"asentiert und bekommt zwei Scores:
\[
  \mathit{prec}_i = \frac{\mathbf{1}^\top(\mathbf{m}_i \wedge \mathbf{y})}{\mathbf{1}^\top \mathbf{m}_i}, \qquad
  \mathit{tp}_i = \mathbf{1}^\top(\mathbf{m}_i \wedge \mathbf{y})
\]

\noindent Falls die Kandidatenanzahl $|\mathcal{R}|$ die Obergrenze \texttt{max\_candidates} \"uberschreitet, greift ein \textbf{Diversity-Filter}: Die Top-$N/2$ nach Precision und die Top-$N/2$ nach TP werden vereinigt. Dadurch bleiben sowohl pr\"azise als auch recall-starke Regeln erhalten.

%% ---------------------------------------------------------------
\section{Phase 2: Greedy-Selektion}

\subsection{Grundidee}

Der Greedy-Algorithmus baut die Composition schrittweise auf: In jedem Schritt wird genau eine Regel ausgew\"ahlt und per OR zur bestehenden Composition hinzugef\"ugt. Das Auswahlkriterium ist die \textbf{marginale Precision} --- die Precision der \emph{gesamten} Composition nach Hinzuf\"ugen der neuen Regel.

\subsection{Formale Beschreibung}

Sei $\mathbf{c}^{(t)}$ der OR-Vektor der Composition nach Schritt~$t$:
\[
  \mathbf{c}^{(0)} = \mathbf{0}, \qquad \mathbf{c}^{(t)} = \mathbf{c}^{(t-1)} \;\vee\; \mathbf{m}_{s_t}
\]
wobei $s_t$ die in Schritt~$t$ gew\"ahlte Regel ist. Definiere:
\begin{align*}
  \mathit{tp}^{(t)} &= \mathbf{1}^\top(\mathbf{c}^{(t)} \wedge \mathbf{y}) &&\text{(kumulative True Positives)}\\
  \mathit{pos}^{(t)} &= \mathbf{1}^\top \mathbf{c}^{(t)} &&\text{(kumulative Positives)}
\end{align*}

\noindent F\"ur eine Kandidatenregel $R_i$ mit Maske $\mathbf{m}_i$ berechne die \emph{neuen, noch nicht abgedeckten} Treffer:
\begin{align*}
  \mathbf{u}_i &= \mathbf{m}_i \;\wedge\; \neg\,\mathbf{c}^{(t-1)} &&\text{(uncovered = neue Treffer durch $R_i$)}\\
  \Delta\mathit{tp}_i &= \mathbf{1}^\top(\mathbf{u}_i \wedge \mathbf{y}) &&\text{(neue True Positives)}\\
  \Delta\mathit{pos}_i &= \mathbf{1}^\top \mathbf{u}_i &&\text{(neue Positives insgesamt)}
\end{align*}

\noindent Die \textbf{Auswahlregel} in Schritt~$t$ ist:
\[
  s_t = \arg\max_{i \in \mathcal{A}^{(t)}} \;\frac{\mathit{tp}^{(t-1)} + \Delta\mathit{tp}_i}{\mathit{pos}^{(t-1)} + \Delta\mathit{pos}_i}
\]
unter der Nebenbedingung $\Delta\mathit{tp}_i \geq \texttt{min\_new\_tp}$.

\subsection{Alive-Menge und Terminierung}

Die Menge der ``lebenden'' Kandidaten $\mathcal{A}^{(t)}$ wird in jedem Schritt aktualisiert:
\begin{itemize}[nosep]
  \item Falls $\Delta\mathit{tp}_i < \texttt{min\_new\_tp}$: Regel $R_i$ wird \textbf{permanent entfernt}. Da die OR-Maske~$\mathbf{c}$ in jedem Schritt nur w\"achst, ist $\Delta\mathit{tp}_i$ monoton fallend --- eine Regel, die jetzt unter dem Schwellwert liegt, kann ihn sp\"ater nie mehr erreichen.
\end{itemize}

\noindent Der Algorithmus terminiert, wenn:
\begin{enumerate}[nosep]
  \item Keine Regel in $\mathcal{A}^{(t)}$ hat $\Delta\mathit{tp}_i \geq \texttt{min\_new\_tp}$, oder
  \item \texttt{max\_steps} erreicht ist.
\end{enumerate}

%% ---------------------------------------------------------------
\section{Algorithmus}

\begin{algorithm}[H]
\caption{Greedy Set Cover}
\begin{algorithmic}[1]
\Require Kandidatenregeln $\{(R_i, \mathbf{m}_i)\}$, Labels $\mathbf{y}$, Parameter \texttt{min\_new\_tp}
\Ensure Schrittfolge $\mathcal{S} = [(t, \mathit{recall}^{(t)}, \mathit{prec}^{(t)}, \text{Regeln}^{(t)})]$
\State $\mathbf{c} \gets \mathbf{0}^n$, $\;\mathit{tp} \gets 0$, $\;\mathit{pos} \gets 0$
\State $\mathcal{A} \gets \{1, \dots, |\mathcal{R}|\}$ \Comment{Alive-Menge}
\For{$t = 1, 2, \dots$}
  \State $\mathit{best\_score} \gets -1$, $\;\mathit{best} \gets \text{nil}$
  \State $\mathcal{A}' \gets \emptyset$
  \For{$i \in \mathcal{A}$}
    \State $\mathbf{u}_i \gets \mathbf{m}_i \wedge \neg\,\mathbf{c}$
    \State $\Delta\mathit{tp}_i \gets \mathbf{1}^\top(\mathbf{u}_i \wedge \mathbf{y})$
    \If{$\Delta\mathit{tp}_i < \texttt{min\_new\_tp}$}
      \State \textbf{continue} \Comment{Permanent entfernt ($\Delta\mathit{tp}$ monoton fallend)}
    \EndIf
    \State $\mathcal{A}' \gets \mathcal{A}' \cup \{i\}$
    \State $\Delta\mathit{pos}_i \gets \mathbf{1}^\top \mathbf{u}_i$
    \State $\mathit{score}_i \gets (\mathit{tp} + \Delta\mathit{tp}_i) \;/\; (\mathit{pos} + \Delta\mathit{pos}_i)$
    \If{$\mathit{score}_i > \mathit{best\_score}$}
      \State $\mathit{best\_score} \gets \mathit{score}_i$, $\;\mathit{best} \gets i$
    \EndIf
  \EndFor
  \State $\mathcal{A} \gets \mathcal{A}'$
  \If{$\mathit{best} = \text{nil}$}
    \State \textbf{break} \Comment{Keine Regel erf\"ullt \texttt{min\_new\_tp}}
  \EndIf
  \State $\mathbf{c} \gets \mathbf{c} \;\vee\; \mathbf{m}_{\mathit{best}}$
  \State $\mathit{tp} \gets \mathit{tp} + \Delta\mathit{tp}_{\mathit{best}}$, $\;\mathit{pos} \gets \mathit{pos} + \Delta\mathit{pos}_{\mathit{best}}$
  \State Speichere Schritt $(t, \;\mathit{tp}/\text{total\_pos}, \;\mathit{tp}/\mathit{pos}, \;\text{Regeln})$
\EndFor
\end{algorithmic}
\end{algorithm}

%% ---------------------------------------------------------------
\section{Beispiel}

Gegeben seien 4 Kandidatenregeln mit folgenden Maskenvektoren ($n = 10$, $\text{total\_pos} = 5$, \texttt{min\_new\_tp} $= 1$):

\medskip
\begin{center}
\renewcommand{\arraystretch}{1.2}
\begin{tabular}{lcccc}
\toprule
& $R_1$ & $R_2$ & $R_3$ & $R_4$ \\
\midrule
Maske $\mathbf{m}_i$ & \texttt{1100100000} & \texttt{0011010000} & \texttt{0000011100} & \texttt{0000000011} \\
TP & 2 & 2 & 1 & 0 \\
Pos & 3 & 3 & 3 & 2 \\
Precision & 66.7\% & 66.7\% & 33.3\% & 0\% \\
\bottomrule
\end{tabular}
\end{center}

\medskip
\noindent Label: $\mathbf{y} = \texttt{1101010000}$ (Positives an Stellen 1, 2, 4, 6).

\bigskip
\textbf{Schritt 1:} Alle Regeln werden evaluiert. $R_4$ hat $\mathit{TP} = 0$ und wird entfernt. $R_1$ und $R_2$ haben marginale Precision $\frac{2}{3} = 66{,}7\%$. $R_1$ wird gew\"ahlt (Tie-Break).
\[
  \mathbf{c}^{(1)} = \texttt{1100100000}, \quad \mathit{tp} = 2, \quad \mathit{pos} = 3, \quad \text{Recall} = \frac{2}{5} = 40\%
\]

\textbf{Schritt 2:} Nur noch uncovered Treffer z\"ahlen. $R_2$: $\mathbf{u}_2 = \texttt{0011010000}$, $\Delta\mathit{tp} = 2$, $\Delta\mathit{pos} = 3$. Marginale Precision: $\frac{2+2}{3+3} = 66{,}7\%$. $R_3$: $\Delta\mathit{tp} = 1$, $\Delta\mathit{pos} = 3$, marginale Precision: $\frac{2+1}{3+3} = 50\%$. $\Rightarrow$ $R_2$ wird gew\"ahlt.
\[
  \mathbf{c}^{(2)} = \texttt{1111110000}, \quad \mathit{tp} = 4, \quad \mathit{pos} = 6, \quad \text{Recall} = 80\%
\]

\textbf{Schritt 3:} $R_3$: $\mathbf{u}_3 = \texttt{0000001100}$, $\Delta\mathit{tp} = 0$ $\Rightarrow$ entfernt. Keine Kandidaten mehr. \textbf{Terminierung.}

\bigskip
\noindent\textbf{Ergebnis:} 2 Regeln, Recall = 80\%, Precision = 66{,}7\%. F\"ur Recall $\geq 40\%$ w\"ahlt man Schritt~1 (1~Regel, 66{,}7\%), f\"ur $\geq 80\%$ Schritt~2 (2~Regeln, 66{,}7\%).

%% ---------------------------------------------------------------
\section{Eigenschaften}

\subsection{Determinismus und Monotonie}

Der Algorithmus ist \textbf{vollst\"andig deterministisch}: Bei gleichen Eingaben und Parametern ergibt sich stets dieselbe Schrittfolge. Da in jedem Schritt mindestens $\texttt{min\_new\_tp} \geq 1$ neue True Positives hinzukommen, ist die Recall-Folge \textbf{streng monoton steigend}:
\[
  \mathit{recall}^{(0)} = 0 \;<\; \mathit{recall}^{(1)} \;<\; \cdots \;<\; \mathit{recall}^{(T)}
\]

\noindent Die Precision-Folge ist dagegen im Allgemeinen \textbf{monoton fallend}, da sp\"atere Regeln typischerweise unpr\"aziser sind.

\subsection{Einmaliges Laufen statt mehrfacher Optimierung}

Da der Algorithmus deterministisch ist, gen\"ugt \textbf{ein einziger Lauf} bis zur Terminierung. Aus der gespeicherten Schrittfolge kann f\"ur jedes beliebige Recall-Ziel $r^*$ die optimale Composition abgelesen werden:
\[
  \text{Schritt}^*(r^*) = \min\{t : \mathit{recall}^{(t)} \geq r^*\}
\]
Dies ist wesentlich effizienter als separate L\"aufe f\"ur verschiedene Recall-Ziele.

\subsection{Rolle von \texttt{min\_new\_tp}}

Der Parameter \texttt{min\_new\_tp} steuert den Trade-off zwischen Pr\"azision und Recall-Abdeckung:

\begin{center}
\renewcommand{\arraystretch}{1.3}
\begin{tabular}{lp{10cm}}
\toprule
\texttt{min\_new\_tp} & Auswirkung \\
\midrule
$= 1$ & Jede Regel mit mindestens einem neuen TP wird ber\"ucksichtigt. Maximaler erreichbarer Recall, aber viele unpr\"azise Regeln in sp\"ateren Schritten. \\
$= 5$ & Balance: gen\"ugend Granularit\"at bei akzeptabler Regelqualit\"at. \\
$= 10$ & Nur Regeln mit substanziellem Beitrag werden gew\"ahlt. Weniger Schritte, aber Recall-L\"ucken m\"oglich. \\
\bottomrule
\end{tabular}
\end{center}

\subsection{Adaptive Schwelle}

In der Variante mit festem Recall-Ziel (\texttt{greedy\_select}) wird \texttt{min\_new\_tp} nicht konstant gehalten, sondern \textbf{adaptiv} an den verbleibenden Bedarf angepasst:
\[
  \texttt{adaptive\_min}^{(t)} = \max\!\Big(\texttt{min\_new\_tp},\;\big\lfloor 0{,}01 \cdot (\texttt{min\_tp} - \mathit{tp}^{(t-1)})\big\rfloor\Big)
\]
wobei $\texttt{min\_tp} = \lceil \texttt{min\_recall} \cdot \text{total\_pos} \rceil$ die Mindestanzahl an TPs ist.

\medskip
\noindent\textbf{Intuition:} Am Anfang fehlen noch viele TPs (z.\,B.\ $\texttt{remaining} = 2{.}800$), sodass $1\%$ davon $= 28$ ist --- nur Regeln mit substanziellem Beitrag werden gew\"ahlt. Gegen Ende ($\texttt{remaining} = 50$) sinkt die Schwelle auf $\max(\texttt{min\_new\_tp}, 0) = \texttt{min\_new\_tp}$, sodass auch kleine Regeln zugelassen werden, um den letzten Recall zu erreichen.

\medskip
\noindent\textbf{Konsequenz f\"ur die Alive-Menge:} Da die adaptive Schwelle sinken kann, d\"urfen Regeln mit $\Delta\mathit{tp}_i < \texttt{adaptive\_min}^{(t)}$ \emph{nicht} permanent entfernt werden --- sie k\"onnten in einem sp\"ateren Schritt mit niedrigerer Schwelle wieder in Frage kommen. Hier wird nur bei $\Delta\mathit{tp}_i < 1$ permanent entfernt (anders als bei der Variante mit fixem \texttt{min\_new\_tp}).

\medskip
\noindent\textbf{Effekt:} Fr\"uhe Schritte w\"ahlen breite, recall-starke Regeln; sp\"ate Schritte f\"ullen mit kleineren, pr\"azisen Regeln auf. Dadurch entsteht ein nat\"urlicher \"Ubergang von grober zu feiner Abdeckung.

%% ---------------------------------------------------------------
\section{Vergleich mit B\&B-Ans\"atzen}

\begin{center}
\renewcommand{\arraystretch}{1.4}
\begin{tabular}{p{3.5cm}p{5.5cm}p{5.5cm}}
\toprule
& \textbf{Greedy Set Cover} & \textbf{B\&B-Ans\"atze} \\
\midrule
Optimalit\"at & Greedy-Approximation (keine Garantie f\"ur globales Optimum) & Exakt optimal (innerhalb des Suchraums) \\
Laufzeit & $O(T \cdot |\mathcal{R}| \cdot n)$ mit $T$ Schritten & Exponentiell im Worst Case \\
Skalierbarkeit & Problemlos f\"ur Millionen Zeilen und Tausende Kandidaten & Begrenzt durch $\binom{R}{r}$ Combinations \\
Recall-Constraint & Implizit durch Schrittfolge & Explizit als Pruning-Kriterium \\
Mehrere Recall-Ziele & Ein Lauf f\"ur alle & Je ein separater Lauf \\
\bottomrule
\end{tabular}
\end{center}

%% ---------------------------------------------------------------
\section{Komplexit\"at}

\renewcommand{\arraystretch}{1.3}
\begin{center}
\begin{tabular}{lll}
\toprule
Phase & Komplexit\"at & Beschreibung \\
\midrule
Apriori & $O(m^k \cdot n)$ & Level-weiser Aufbau, $m$ Pr\"adikate, $k$ Levels \\
Scoring & $O(|\mathcal{R}| \cdot n)$ & Einmalige Precision/TP-Berechnung \\
Greedy & $O(T \cdot |\mathcal{R}| \cdot n)$ & $T$ Schritte, je $|\mathcal{R}|$ Kandidaten \\
\bottomrule
\end{tabular}
\end{center}

\noindent Praktisch: Die Alive-Menge $\mathcal{A}$ schrumpft in jedem Schritt, sodass sp\"atere Schritte schneller sind. F\"ur $k=3$, $|\mathcal{R}| \approx 3{.}000$ Kandidaten und $n = 2{,}8$M Zeilen dauert ein vollst\"andiger Lauf mit $\sim$70 Schritten wenige Sekunden.

%% ---------------------------------------------------------------
\section{Nachverarbeitung}

Nach der Greedy-Selektion enth\"alt die Composition m\"oglicherweise redundante Regeln. Drei Nachverarbeitungsschritte reduzieren die Regelanzahl, ohne den Recall zu verringern (bei Subsumption und Unique-TP) bzw.\ mit kontrolliertem Recall-Verlust (bei Merge).

\subsection{Subsumption Removal}

\subsubsection{Threshold-Hierarchie}

Pr\"adikate derselben Feature-Gruppe bilden eine \textbf{Lockerheitshierarchie}. Strenge Schwellwerte matchen weniger Zeilen als lockere:
\[
  \texttt{eq} \;\subset\; \texttt{dist\_1} \;\subset\; \texttt{dist\_2} \;\subset\; \texttt{dist\_3} \;\subset\; \texttt{dist\_4}
\]
\[
  \texttt{prop\_0.99} \;\subset\; \texttt{prop\_0.95} \;\subset\; \texttt{prop\_0.9} \;\subset\; \texttt{prop\_0.8} \;\subset\; \texttt{prop\_0.7}
\]
\[
  \texttt{overlap\_0.75} \;\subset\; \texttt{overlap\_0.5} \;\subset\; \texttt{overlap\_0}
\]
Jeder Schwellwert $\tau$ hat einen \textbf{Lockerheitsindex} $\ell(\tau) \in \mathbb{N}_0$, wobei $\ell = 0$ der strengste ist. F\"ur zwei Schwellwerte $\tau_1, \tau_2$ derselben Familie mit $\ell(\tau_1) \leq \ell(\tau_2)$ gilt:
\[
  \mathbf{m}(\tau_1\_f) \;\subseteq\; \mathbf{m}(\tau_2\_f)
\]
d.\,h.\ jede Zeile, die den strengen Schwellwert erf\"ullt, erf\"ullt auch den lockeren.

\subsubsection{Dominanz-Relation}

Eine Regel $A$ \textbf{dominiert} eine Regel $B$ (geschrieben $A \succ B$), wenn:
\begin{enumerate}[nosep]
  \item F\"ur jedes Pr\"adikat $p_A$ in $A$ existiert ein Pr\"adikat $p_B$ in $B$ aus \emph{derselben} Feature-Gruppe mit $\ell(p_A) \geq \ell(p_B)$ (d.\,h.\ $A$ ist mindestens so locker wie $B$ in dieser Gruppe).
  \item $B$ darf zus\"atzliche Pr\"adikate besitzen (die $B$ noch enger machen).
\end{enumerate}

\noindent Falls $A \succ B$, dann gilt f\"ur die Masken: $\mathbf{m}_B \subseteq \mathbf{m}_A$. Jede Zeile, die von $B$ abgedeckt wird, wird auch von $A$ abgedeckt. In einer OR-Composition ist $B$ daher \textbf{redundant}.

\subsubsection{Beispiel}

Gegeben seien drei ausgew\"ahlte Regeln:
\begin{align*}
  R_1 &= \{\texttt{dist\_2\_end\_date}\} \\
  R_2 &= \{\texttt{eq\_end\_date},\; \texttt{eq\_start\_date}\} \\
  R_3 &= \{\texttt{dist\_1\_end\_date},\; \texttt{overlap\_0.5\_X}\}
\end{align*}

\noindent Pr\"ufe $R_1 \succ R_2$: $R_1$ hat \texttt{dist\_2\_end\_date} (Gruppe \texttt{end\_date}, $\ell = 2$). $R_2$ hat \texttt{eq\_end\_date} in derselben Gruppe mit $\ell = 0 \leq 2$ \checkmark. $R_2$ hat ein zus\"atzliches Pr\"adikat \texttt{eq\_start\_date} \checkmark. $\Rightarrow$ $R_1 \succ R_2$: \textbf{$R_2$ wird entfernt.}

\medskip
\noindent Pr\"ufe $R_1 \succ R_3$: $R_1$ hat \texttt{dist\_2\_end\_date} ($\ell = 2$). $R_3$ hat \texttt{dist\_1\_end\_date} ($\ell = 1 \leq 2$) \checkmark. $R_3$ hat \texttt{overlap\_0.5\_X} zus\"atzlich \checkmark. $\Rightarrow$ $R_1 \succ R_3$: \textbf{$R_3$ wird entfernt.}

\medskip
\noindent\textbf{Ergebnis:} Von 3 Regeln bleibt nur $R_1$ \"ubrig. Der Recall bleibt unver\"andert, da $\mathbf{m}_{R_2} \cup \mathbf{m}_{R_3} \subseteq \mathbf{m}_{R_1}$. Die Precision kann nur steigen (weniger FP durch weniger Regeln).

\subsubsection{Algorithmus}

F\"ur jedes Regelpaar $(R_i, R_j)$ mit $|R_i| \leq |R_j|$ wird gepr\"uft, ob $R_i \succ R_j$. Falls ja, wird $R_j$ als dominiert markiert. Komplexit\"at: $O(r^2 \cdot k)$ mit $r$ Regeln und maximal $k$ Pr\"adikaten pro Regel.

%% ---------------------------------------------------------------
\subsection{Unique-TP Removal}

Nach der Subsumption-Bereinigung kann es Regeln geben, die zwar nicht von einer einzelnen anderen Regel dominiert werden, deren True Positives aber \emph{insgesamt} bereits von der Vereinigung aller anderen Regeln abgedeckt sind.

\subsubsection{Definition}

F\"ur eine Regel $R_i$ in der Composition $\{R_1, \dots, R_r\}$ sind die \textbf{Unique True Positives}:
\[
  \mathit{UTP}(R_i) = \mathbf{1}^\top\!\left(\mathbf{m}_i \;\wedge\; \neg\!\bigvee_{j \neq i} \mathbf{m}_j \;\wedge\; \mathbf{y}\right)
\]
Das sind die TPs, die \emph{ausschlie\ss lich} von $R_i$ abgedeckt werden und von keiner anderen Regel in der Composition.

\medskip
\noindent Falls $\mathit{UTP}(R_i) = 0$: Entfernung von $R_i$ \"andert den Recall nicht (alle TPs von $R_i$ werden von anderen Regeln abgedeckt), aber die Precision steigt, da die FPs von $R_i$ wegfallen.

\subsubsection{Beispiel}

\begin{center}
\renewcommand{\arraystretch}{1.2}
\begin{tabular}{lccc}
\toprule
& Maske & TP-Maske & UTP \\
\midrule
$R_1$ & \texttt{111000} & \texttt{110000} & 1 (Pos.~1) \\
$R_2$ & \texttt{011100} & \texttt{010100} & 1 (Pos.~4) \\
$R_3$ & \texttt{010010} & \texttt{010000} & 0 \\
\bottomrule
\end{tabular}
\end{center}

\noindent $\mathbf{y} = \texttt{110100}$. $R_3$ hat TP an Position~2, aber $R_1$ und $R_2$ decken Position~2 ebenfalls ab. $\Rightarrow$ $\mathit{UTP}(R_3) = 0$, $R_3$ wird entfernt. Recall bleibt bei $\frac{3}{3} = 100\%$, aber die FP von $R_3$ (Position~5) f\"allt weg.

\subsubsection{Algorithmus}

F\"ur jede Regel $R_i$: Berechne die OR-Maske aller anderen Regeln $\bigvee_{j \neq i} \mathbf{m}_j$, dann $\mathit{UTP}(R_i)$. Falls $\mathit{UTP}(R_i) = 0$, entferne $R_i$. Komplexit\"at: $O(r^2 \cdot n)$.

%% ---------------------------------------------------------------
\subsection{Merge}

Der Merge-Schritt reduziert die Anzahl der Regeln durch \textbf{Zusammenfassung \"ahnlicher Regeln}. Regeln, die viele gemeinsame Pr\"adikate teilen, werden zu einer einzigen, strengeren Regel vereinigt.

\subsubsection{Vorgehen}

\begin{enumerate}[nosep]
  \item \textbf{Paar-H\"aufigkeiten z\"ahlen:} F\"ur jedes Pr\"adikatpaar $(p_a, p_b)$ wird gez\"ahlt, in wie vielen Regeln beide gleichzeitig vorkommen.
  \item \textbf{Greedy-Gruppierung:} Das h\"aufigste Paar wird ausgew\"ahlt. Alle Regeln, die dieses Paar enthalten, bilden eine Gruppe.
  \item \textbf{Vereinigung:} Alle Pr\"adikate der Gruppe werden zu einer einzigen AND-Regel zusammengefasst (Mengenvereinigung).
  \item Wiederhole f\"ur das n\"achsth\"aufigste Paar (nur mit noch nicht gruppierten Regeln).
\end{enumerate}

\subsubsection{Beispiel}

Gegeben seien 4 Regeln:
\begin{align*}
  R_1 &= \{A,\; B,\; C\} \\
  R_2 &= \{A,\; B,\; D\} \\
  R_3 &= \{A,\; B,\; E\} \\
  R_4 &= \{C,\; D,\; F\}
\end{align*}

\noindent Das Paar $(A, B)$ kommt in $R_1, R_2, R_3$ vor (3-mal). Diese drei Regeln werden zusammengefasst:
\[
  R_1 \cup R_2 \cup R_3 = \{A,\; B,\; C,\; D,\; E\}
\]
$R_4$ bleibt unver\"andert. \textbf{Ergebnis:} 2 Regeln statt 4.

\subsubsection{Auswirkung auf Metriken}

Der Merge ist \textbf{nicht recall-erhaltend}. Die urspr\"ungliche Composition
\[
  (A \land B \land C) \;\lor\; (A \land B \land D) \;\lor\; (A \land B \land E)
\]
wird zu
\[
  A \land B \land C \land D \land E
\]
Die zusammengefasste Regel verlangt \emph{alle} Pr\"adikate gleichzeitig (AND), w\"ahrend die urspr\"unglichen Regeln per OR verkn\"upft waren. Dadurch matcht die Merge-Regel \emph{weniger} Zeilen:
\[
  \mathbf{m}_{A \land B \land C \land D \land E} \;\subseteq\; \mathbf{m}_{A \land B \land C} \;\cup\; \mathbf{m}_{A \land B \land D} \;\cup\; \mathbf{m}_{A \land B \land E}
\]

\noindent Der Vorteil liegt in der \textbf{Interpretierbarkeit}: Weniger, aber l\"angere Regeln sind oft leichter zu verstehen als viele kurze. Au\ss erdem steigt typischerweise die Precision, da die Merge-Regel strenger filtert.

\medskip
\noindent\textbf{Anwendung:} Der Merge wird nach der Greedy-Selektion und den recall-erhaltenden Bereinigungen (Subsumption, Unique-TP) durchgef\"uhrt. Recall und Precision werden danach neu berechnet.

\end{document}
